\section{Question and contribution}
Our Research Under Uncertainty (RUU) program studies how models form assumptions, choose
experiments and revise predictions when an environment's mechanism is unknown.
Decision records offer a way to examine those activities. Before using them
as evidence, however, we need to know which relationships they measure.
This paper asks: \textbf{when does a record that matches a model's action
identify a selection rule, and when does it help predict a later choice?}

We study externally supplied tables and descriptions in two simulated task
families. An executable record combines an action-scoring representation
with a deterministic rule for selecting an action. For an incompletely
specified record, the analyst must supply the remaining choices; we call
this a \emph{selector completion}. Combining one record's representation
with another record's rule is a \emph{crossing}. For example, applying
column means and column maxima to the same table can select different
actions. Our initial component diagnostic used such crossings to examine
which part of a record supported a successful choice.

\paragraph{A concrete identification problem.}
In the pilot's original representation A2, one plan is best on all five
examined criteria. Total cost and trajectory peak therefore reproduce the
same model action. In a second representation A1, held fixed for comparison,
the criteria separate: total cost selects a plan with a preservation witness,
while trajectory peak selects a different plan with a conditional loss
certificate. Agreement on A2 leaves these completions indistinguishable,
although they give different crossed diagnoses
(Section~\ref{sec:selector-sensitivity}). This example supplies a practical
check: \emph{does the validation situation give candidate selectors an
opportunity to disagree?} When it does not, the diagnosis must remain
conditional on the independently specified selection contract.

\paragraph{Correspondence and prediction are separately measurable.}
A follow-up separates requiring a record from externally selecting its
action. Each profile supplies eight records matching its own executed
proposal. Fourteen of these sixteen numeric tables have one row, so both
supported rules necessarily agree; their agreement in the remaining two
multirow tables is observed. We pair these counts with an action-copy control
and a scan over the supported rules (Section~\ref{sec:followup}). This makes
both correspondence and its opportunity for rule discrimination visible.

A prospective pilot then tests added prediction on new choices. Records are
frozen before two target branches per source, and paired observers receive
the same information with or without the record. The observed net accuracy
differences are $-1/8$ for Astra and $0/8$ for Sol across four purposeful
histories per profile (Section~\ref{sec:prospective}). The pilot measures
transfer from a source that reports an action and record to a target that
reports an action alone; it evaluates the numeric table and accompanying
prose together.

\paragraph{Contribution to evaluation practice.}
The evaluation decision at issue is accepting an analyst-completed selector
because it reproduces the source action. The archived A2 case exposes the
problem: one plan is best on all five criteria, so every tested completion
passes that check, yet the same completions produce different choices and
conditional task diagnoses on A1. Researchers using such records should
therefore retain the compatible selector set and test a situation where its
members can disagree before attributing a diagnosis to a particular rule.

The follow-up and prospective panels make two further reporting decisions
concrete: separate a record's match to an unforced proposal from the task
label, and measure added prediction on new choices after freezing the record.
An action change can preserve its P or L label
(Table~\ref{tab:followup-rule-examples}); an aggregate prediction difference
can combine opposite observed contrasts across tasks
(Section~\ref{sec:prospective}). These cases give researchers specific checks
to add to a record evaluation. The common-menu recovery identity and
single-row rule agreement explain what the contract already guarantees
(Appendix~\ref{app:structural-control}); the archived cases and model decisions
show where the distinctions matter in practice. Table~\ref{tab:study-map}
maps each panel to its question, and Section~\ref{sec:related} relates these
external-record measurements to explanation evaluation.

\begin{table}[htbp]
\centering\small
\begin{tabular}{@{}>{\raggedright\arraybackslash}p{0.14\linewidth}>{\raggedright\arraybackslash}p{0.23\linewidth}>{\raggedright\arraybackslash}p{0.25\linewidth}>{\raggedright\arraybackslash}p{0.29\linewidth}@{}}\toprule
Study & Question & Intervention or check & Supported claim\\\midrule
Pilot (8 trials) & Does native agreement identify a selector? & Five analyst completions of archived records; component crossings & Concrete non-identifiability under these completions\\
Transfer (48 trials) & Are free-form records crossing-ready? & Optional records across two task families & Applicability is limited; zero eligible pairs is not success\\
Common menu (32 trials) & What does the fixed numeric contract add? & Optional records versus required records with external selection & Structural recoverability; record requirement and execution remain combined\\
Follow-up (48 assignments) & Can record and proposal match without forced execution? & A optional/model, B required/model, C required/selector & Profile-specific correspondence, record structure and conditional labels; no conflicting C proposal observed\\
Prospective choice prediction (16 targets) & Does a frozen record add prediction value? & Paired observers with/without record; same public evidence and profile & No positive net accuracy gain observed; four histories/profile, two linked targets/history\\\bottomrule
\end{tabular}
\caption{Study map. Counts refer to distinct panels, not pooled replications.
The follow-up has 46 completed assignments and two retained technical losses.
The prospective pilot has 15 original targets plus one separately completed
target; its original interrupted request remains recorded.
The separate finite-class companion is not needed for the main argument.}
\label{tab:study-map}
\end{table}

\section{Task, record, execution and evaluation}
\label{sec:measurement}
\paragraph{Task.}
The actor receives an action/observation history, a remaining objective and a
finite menu, without the hidden mechanism. In the continuous task it chooses
a bounded two-dimensional control to preserve the possibility of reaching
and maintaining a target region under a path constraint. In the route task
it chooses a legal discrete action under remaining route constraints.
The common-menu and prospective panels offer 13 continuous controls or
three route actions, and allow actors up to three analyses of the observed
history. Both families supply the past actions and observed states, the
remaining objective and constraints; the hidden transition mechanism is
withheld. The prospective panel's concrete source--target changes appear in
Table~\ref{tab:prospective-branches}.

\paragraph{Record.}
The actor can submit a numeric table of candidate failure forecasts and a
declared rule. Columns identify menu actions; rows contain values in $[0,1]$.
Uniform mean minimizes the column mean; worst case minimizes the column
maximum; ties use the lowest index. The executable contract specifies action selection; forecast calibration
and the semantic adequacy of rows require separate evaluation.
Table~\ref{tab:real-record} shows one complete numeric table from the follow-up,
before any formal crossing notation is needed.

% Exact numeric values from finals/trial-003.json; no rounding.
\begin{table}[htbp]
\centering\small
\begin{tabular}{lrrr}\toprule
Archived candidate-loss row & Menu action 0 & Menu action 1 & Menu action 2\\\midrule
Only row & 0.35 & 0.30 & 0.25\\\bottomrule
\end{tabular}
\caption{A real B record: Astra, route history \texttt{probe-03}, repeat 1,
trial 3. The declared mean rule selects action 2, as does worst case.
The model proposed action 2 and B executed that proposal. Its bound menu label
is P. The numbers are uncalibrated forecasts; their ordering and the matching
proposal are reproducible, but their semantic adequacy or causal role is not
established. Exact final JSON and identity bindings are in the bundle.}
\label{tab:real-record}
\end{table}


\paragraph{Execution.}
The follow-up separates record requirement from action source
(Table~\ref{tab:arms-early}). In B the model proposal is executed even if a
valid record would select a different action. In C the external selector
determines execution, while the proposal remains observable. Thus B's
record--proposal agreement is observed; C's record--execution agreement is
guaranteed. Exact instructions and schema are in the bundle.

\begin{table}[htbp]
\centering\small
\begin{tabular}{lp{0.22\linewidth}p{0.23\linewidth}p{0.34\linewidth}}\toprule
Arm & Record requirement & Executed action & What agreement can establish\\\midrule
A & Optional & Model proposal & Correspondence, if a valid record exists\\
B & Required valid record & Model proposal & Observed record--proposal correspondence\\
C & Required valid record & Declared record selector & Execution match by construction; proposal match remains observable\\\bottomrule
\end{tabular}
\caption{Follow-up conditions. A valid record need not match the proposal to
be accepted in B. A missing record leaves correspondence undefined.}
\label{tab:arms-early}
\end{table}

\paragraph{Evaluation.}
The analyst attaches a task label to the executed menu action: preservation
(P), loss (L), or unresolved (U), using the bound feasibility evaluation.
That label is separate from whether a record is valid, matches the proposal,
or changes action under rule substitution. The model does not receive the
privileged evaluator during collection.

\subsection{Why matching needs a control}
For any proposed index $a$ in an $n$-action menu, construct one row
$A_{1j}=0$ when $j=a$, and $A_{1j}=1$ otherwise. Both supported rules select
$a$. This \emph{action-copy negative control} requires only the menu size
and proposed index, with no history, mechanism or task label. Conditional on
the proposal it supplies no independent information about how that proposal
was produced. It nevertheless satisfies validity and correspondence checks.
Consequently those checks alone cannot establish an informative explanation.
The construction tests the sufficiency of those checks. How the participant
produced its record and whether its forecasts are informative remain empirical
questions.

A positive arithmetic control has two rows,
$A=\left(\begin{smallmatrix}0&0.4\\0.7&0.4\end{smallmatrix}\right)$.
Mean scores $(0.35,0.4)$ select action 0; worst-case scores $(0.7,0.4)$
select action 1. The same reader therefore detects a rule contrast when the
values support one. Neither constructed control has a physical task label
or constitutes an additional model trial. The distributed check tests the
copy construction at every proposal index for menu sizes 3 and 13, under
each declared rule, and tests the positive control under both rules.

\subsection{Two external comparisons}
Let $A_i$ be record $i$'s table over a common finite menu and let $r_i$
be its declared deterministic rule. The action $a_i=r_i(A_i)$ is one
column of that menu. Here, $a_i$ denotes the record-selected action; the
model-proposed and executed actions are separate quantities. The
\emph{observed-rule crossing} has four entries
$c_{ij}=r_j(A_i)$ for $i,j\in\{1,2\}$. The separate
\emph{supported-rule-family table} evaluates each $A_i$ under both allowed
rules: uniform mean and worst case, with lowest-index tie breaking.
The family table may contain a rule that neither paired actor used. Its
entries are deterministic evaluations of archived records, with task labels
from the bound menu census; a live model counterfactual would require a new
model decision.

\subsection{Feasibility labels and access}
P requires a successful admissible continuation witness after the action.
L requires the specified loss certificate; absence of a witness alone is not
L. U preserves unresolved numerical cases. Continuous-control labels are
conditional on the disclosed floating-point and bounding premises in the
supplement. The discrete route family's finite computation has a different
verification contract. Both actors and analyst see the legal histories; only
the analyst receives the hidden mechanisms and privileged evaluator during
the original experiment. Reviewer disclosure now permits inspection of these
mechanisms without retroactively changing what participants received.

Record validity, correspondence and rule sensitivity can be checked without
the physical label map. P/L/U comparisons add its assumptions. The exact
arithmetic helper does not prove the uniform library bounds or attest the
historical runtime; those premises remain open. These are distinct evidence
levels when interpreting outcome differences.

\section{Evidence sets and their denominators}
The first three evidence sets below use crossing-pair denominators.
The separate follow-up in Section~\ref{sec:followup} instead compares A/B/C
conditions within eight matched triplets per profile; these condition pairs
are not crossing pairs. Its 48 assignments include two technical losses.
\begin{table}[htbp]
\centering\small
\begin{tabular}{p{0.28\linewidth}rrrr}
\toprule
Evidence set & Trials & Pairs & Eligible records & Eligible pairs\\
\midrule
Pilot (component study) & 8 & 4 & \multicolumn{2}{c}{2 selected component pairs}\\
Prospective transfer & 48 & 24 & 1 & 0\\
Optional record & 16 & 8 & 1 & 0\\
Mandatory record + execution & 16 & 8 & 16 & 8\\
\bottomrule
\end{tabular}
\caption{The three predecessor evidence sets. Pilot component eligibility is a different
construction from the frozen prospective adapter and common-menu contract.
Empty pair denominators leave the crossing effect unassessed.}
\end{table}

The pilot used four fixed histories with two repeats under the requested
gpt-6-astra/max profile. Its component study
focused on two boundary histories, pilot contexts B and D, after outcomes
were available. These history names are unrelated to the follow-up arms A/B/C.
The transfer panel used 48 live trials: two families, three instances per
family, two histories per instance, the gpt-6-astra/max and
gpt-6.1-sol/max profiles, and two repeats. Under its frozen adapter, a record had to
support the prescribed operation without post hoc completion. Only one trial
qualified, leaving zero of 24 pairs crossable. This measures the adapter's applicability to the available external
records. With no eligible pair, the crossing effect is unassessed; internal
representations are outside the adapter's observations. Native transfer labels were 42 P and 6 L. The adapter's mutually exclusive
outcomes were: 24 route trials with no analysis attempt, 18 whose last successful
analysis lacked a table, three whose rule did not reproduce the action, two
with a menu-identity mismatch, and one supported record. These reasons
distinguish absent evidence from an invalid submitted record.
The other 24, continuous-task trials made 72 analysis attempts: 67 succeeded
and five returned errors. The accounting distinguishes an unused analysis opportunity from a failed
analysis call. It measures visible tool use, leaving internal reasoning
outside the observation.

The 32-trial panel crossed two task families, two purposefully selected
histories per family, two model profiles (gpt-6-astra/max and gpt-6.1-sol/max),
two repeats and two conditions. Each condition has 16 trials and eight pairs.
A crossing pair comprises the two repeats at the same history and profile
within one condition, not a pair of trials from different conditions.
Both conditions used a common finite menu and structured final output.
Appendix~\ref{app:common-menu-protocol} identifies the selected histories,
execution contract, resource limits and recorded usage.
The optional condition permitted an observational record; the mandatory
condition required a record and let an external deterministic executor select
the action. Mandatory record--action agreement is consequently a construction
property, unlike a matching optional observational record. The combined
intervention does not isolate a pure reporting effect.

All 32 trials were evaluated: no unstarted slots, technical missingness or
protocol-invalid trials remain. This does not mean the analysis interface was
error-free: 48 attempts yielded 41 successes and seven code errors in seven
trials (five optional, two mandatory). Five errors used attribute access on
the dictionary history; two used the unavailable built-in \texttt{next}.
The prompt described \texttt{history.states}/\texttt{history.actions}, although
the worker passed a dictionary, and did not enumerate the built-ins.
Each affected trial consumed all three analysis slots, one on the error and
two successfully. Appendix~\ref{app:interface-audit} gives the trial-level
accounting with all original trials and labels retained. The effect on
subsequent actions and record structure is unassessed.
Optional P/L/U counts were 13/3/0; mandatory
counts were 10/5/1. Four purposefully chosen histories do not form a population
sample of tasks or support a model ranking. Table~\ref{tab:label-contracts}
separates the families' label contracts. The study inventory in
Appendix~\ref{app:inventory} identifies excluded historical and development
work; the later acquisition experiment has its own denominator.

\begin{table}[htbp]
\centering\small
\begin{tabular}{llrrrl}\toprule
Condition & Family & P & L & U & Verification contract\\\midrule
Optional & Continuous & 5 & 3 & 0 & Numerical, conditional\\
Mandatory & Continuous & 2 & 5 & 1 & Numerical, conditional\\
Optional & Route & 8 & 0 & 0 & Finite integer recurrence\\
Mandatory & Route & 8 & 0 & 0 & Finite integer recurrence\\\bottomrule
\end{tabular}
\caption{Native outcomes in the common-menu panel, eight trials per row.
Continuous P uses an admissible continuation witness; continuous L depends
on the numerical and historical-runtime premises in Appendix F. U remains
unresolved. Successful present-host replay does not close those premises.}
\label{tab:label-contracts}
\end{table}

\section{Results: what agreement and crossing identify}
\subsection{Pilot counterexample: native agreement does not identify the selector}
\label{sec:selector-sensitivity}
The missing choice can be seen in pilot context B. Its two archived model
analyses supplied different candidate dynamics and planning ingredients;
the native-loss representation A1 has eight candidates, while the
native-preservation representation A2 has four. The researcher reconstruction
computes one multi-step plan per candidate. To execute this component on
either representation, an analyst must also specify \emph{which plan supplies
the next action}. The declared completion R2* chooses the plan with the
smallest total penalized cost. This last selection is an analyst choice
(Appendix~\ref{app:components}), not identified by reproducing the native action.

Here \emph{native-action-compatible} means only that the selected first action
on A2 is within $10^{-6}$ of the recorded native action. It does not mean
agreement with every semantic commitment of the model's narrative or recovery
of its internal policy. We examine five selectors over the same fixed plans:
total cost, maximum target-step norm, mean squared target-step norm, trajectory
peak norm and final-step norm. These are explicit summaries of the plans'
objective, target attainment or path behavior, making the analyst's choice
inspectable. They obey the same action bounds and lowest-index tie breaking,
but need not encode the same preferences; final-step norm alone, for example,
is not the full multi-step task criterion. This is a finite retrospective
sensitivity family, not an exhaustive set of faithful completions.

\begin{table}[htbp]
\centering\small
\begin{tabular}{lccr}\toprule
Criterion minimized & Native A2: plan/label & Crossed A1: plan/label & Action norm\\\midrule
Total penalized cost (R2*) & 2/P & 2/P & 0.575465\\
Maximum target-step norm & 2/P & 2/P & 0.575465\\
Mean squared target-step norm & 2/P & 2/P & 0.575465\\
Trajectory peak norm & 2/P & 3/L & 0.042684\\
Final-step norm & 2/P & 5/P & 0.579114\\
\bottomrule\end{tabular}
\caption{Pilot context B: five selectors applied to the same fixed plan pools.
Indices are zero-based. All five select native A2 plan 2 and reproduce the
original action within $10^{-6}$. The first three select the same crossed
A1 plan; the minimum-peak selector instead selects plan 3, giving L rather
than P. Norms refer to crossed actions. Labels retain the conditional
continuous evaluation premises. These five selectors are not independent
replications or a probability sample of completions.}
\label{tab:b-completions}
\end{table}


All five select A2 plan 2, whose first action matches the native action to
approximately $1.81\times10^{-9}$. Plan 2 has a strictly lower value than
each of the other three plans on every one of the five criteria. A2 thus
provides no opportunity to discriminate among these selectors: agreement
there follows from a common best plan. Now hold the crossed A1 representation,
history and computed plans fixed. Total cost selects A1 plan 2, with action
norm $0.575465$ and a preservation witness. Trajectory peak instead selects
A1 plan 3, with norm $0.042684$ and a conditional loss certificate
(Table~\ref{tab:b-completions}). The reader command recomputes all twelve plans
and five selections, re-evaluates their first actions with the bundled grader,
and compares the complete archived witness/certificate records
(Section~\ref{sec:reader-access}). It therefore checks the label change as well
as the selector arithmetic. The numerical and historical-runtime premises
in Appendix~\ref{app:runtime-proof} remain open.

The resulting component diagnosis depends on which native-compatible
selector the analyst supplies. To use native-action agreement as a validation
test, report the candidate selectors and whether the validation cases
actually separate their choices. An independently specified selection
contract can restrict the alternatives. In this case, the common best plan
on A2 leaves that choice open. The two A1 plans also differ in action
magnitude, so this contrast evaluates the full selection change rather than
isolating a horizon or direction effect.

Under the originally declared completions, pilot context B had two P
off-diagonals and context D two L off-diagonals. Each native pair already
contained one P and one L. Those patterns describe completion-conditional
component dependence, not new outcome categories or internal causal ground
truth. The pilot's heterogeneous optimizers differ from the later common-menu
contract; the latter's structural identity is not extended to this setting.

\subsection{Common-menu applicability and recoverability}
The mandatory condition made all eight pairs executable. It did not make
their crossing tables informative by construction. Every pair used the same
declared rule, so all 32 observed-rule cells were already determined by native
actions and rule identity. Among the 16 mandatory records, 13 had one row,
one contained duplicate rows alongside distinct rows, and 13 had only one
distinct row. A duplicate-containing table can also contain distinct rows, so the
duplicate and single-distinct-row counts measure different properties.

Native actions plus the single-distinct-row check determine 29 of the
32 supported-rule-family cells. The remaining sensitivity can be calculated on each record independently.
Appendix~\ref{app:structural-control} states the fixed-contract identity and
the complete recoverability table. These calculations measure recoverable
cells; improved diagnostic accuracy requires an external diagnostic target.


Fifteen of 16 mandatory records selected the same action under both supported
rules. The remaining record changed action and its conditional continuous
P/L category; the independent
scan recovered this change. One pair involved U and remains unresolved rather
than being converted into a resolved loss. The two 32-cell displays overlap;
after within-pair deduplication they reference 30 unique numerical computations,
not 64 new samples. Cross-action labels reuse the existing bound menu census.
No new live model counterfactuals were collected for this retrospective analysis.

\paragraph{Worked same-record contrast.}
For the mandatory \texttt{probe-01} record from gpt-6-astra, repeat 1,
the independent scan yields:
\begin{center}
\begin{tabular}{lrl}\toprule
Rule applied to this record & Menu index & Bound task label\\\midrule
Uniform mean & 7 & L (conditional continuous certificate)\\
Worst case & 12 & P (continuation witness)\\\bottomrule
\end{tabular}
\end{center}
Indices are zero-based. The recorded table, menu identity and label census
are held fixed. No partner record is needed to expose this sensitivity;
it concerns the declared external rule, not a demonstrated internal cause.
The alternative action's label is a deterministic menu evaluation, not an
additional live model decision. This is the sole rule-sensitive mandatory
record, so it must not be presented as a prevalent population effect.
This trial was also affected by the dictionary-access error.
The archived table's rule sensitivity is reproducible; whether correcting
the interface would have changed that table or the decision is unknown.

The optional condition contains one submitted record, which is eligible,
and 15 trials with no submitted record (\texttt{no\_record}). It has no
eligible pairs. The fifteen absences and one eligible submission have different statuses.
With no eligible pair, the optional crossing effect remains unassessed.
The panel thus locates an applicability bottleneck in the optional condition
and a redundancy bottleneck under mandatory recording with external selection.

\section{Follow-up: required records with and without external execution}
\label{sec:followup}
The earlier mandatory condition combined a record requirement with an external
executor. The follow-up adds a required-record condition that executes the
model's proposal. Each matched triplet uses one history, repeat and
model/max/Fast profile: A permits an optional record and executes the proposal;
B requires a valid record and executes the proposal; C requires a valid record
and executes its declared selector, retaining the proposal for observation.
Every trial starts a fresh context. Four purposefully selected historical
contexts, two profiles and two repeats give 48 planned assignments, eight
triplets per profile. These repeats are not independent draws of new worlds.
The allocation and six possible within-triplet orders were fixed before
collection; earlier-panel trials are not reused as controls.
Appendix~\ref{app:followup} gives the execution segments, interface and reader
boundaries. No new population sample or live crossed decision was collected.

\begin{table}[htbp]
\centering\small
\begin{tabular}{llrrrrrrr}\toprule
Profile & Arm & Done & Tech. & P/L/U & Valid & Match & One row & Same action\\\midrule
Astra & A & 7 & 1 & 6/1/0 & 0 & -- & 0 & 0\\
Astra & B & 8 & 0 & 8/0/0 & 8 & 8/8 & 8 & 8\\
Astra & C & 7 & 1 & 5/2/0 & 7 & 7/7 & 6 & 6\\\midrule
Sol & A & 8 & 0 & 7/1/0 & 0 & -- & 0 & 0\\
Sol & B & 8 & 0 & 6/2/0 & 8 & 8/8 & 6 & 8\\
Sol & C & 8 & 0 & 7/1/0 & 8 & 8/8 & 7 & 7\\\bottomrule
\end{tabular}
\caption{Follow-up, eight planned assignments per row; profiles are separate
strata, not ranked. Done means a completed trial, not full task success.
Valid counts usable records; Match is record-selected action versus model
proposal among valid records. The empty A denominator is undefined.
Same action counts records for which both supported rules select the same
action. Sol C has one duplicate-containing record; all other rows have zero.
There are no protocol losses, active or unstarted assignments.}
\label{tab:followup-records}
\end{table}

\paragraph{Availability and correspondence.}
In A, no valid record was available among Astra's seven or Sol's eight
completed trials. This is absence of the specified external record, not
absence of reasoning. In B, each profile supplied eight valid records and all
eight matched its own proposal. Unlike C's executed-action agreement, this
correspondence was not imposed by the executor. It does not establish that
the record caused the proposal or faithfully represents internal reasoning.
Table~\ref{tab:B-structure} separates the observed record structures from
their mathematical consequences. Astra supplied eight single-row B records;
Sol supplied six single-row and two multirow B records. For every single-row
record, equality of mean and worst case is forced, so it is not an independent
empirical success of rule agreement. Sol's multirow records contain three and
four distinct rows; both also select the same action under the two rules.
That equality is an observation about these values, not forced by having one
distinct row. Their source orders are 44 and 47, respectively. No multirow
conclusion is available for Astra's empty denominator.

% Generated from the preserved minimal finals by contribution_checks.py.
\begin{table}[htbp]
\centering\small
\begin{tabular}{lrrrr}\toprule
Profile & Proposal matches & One row & Multiple rows & Multirow same action\\\midrule
Astra & 8/8 & 8 & 0 & --\\
Sol & 8/8 & 6 & 2 & 2/2\\
\bottomrule\end{tabular}
\caption{B records, separated by profile. One-row mean/worst-case equality is
forced by the contract. For multirow records, equality is an observed property
of the archived values. Astra has no multirow denominator, so its dash is
undefined, not a success. Sol's two multirow records have three and four
distinct rows, respectively.}
\label{tab:B-structure}
\end{table}


The empirical result is the co-occurrence of proposal matching and these
record structures without external selection. The action-copy control in
Section~\ref{sec:measurement} establishes why matching alone leaves
substantive use of the history unresolved. The actual records' predictive
value requires its own measurement. Single-row and duplicate records satisfy
the protocol and remain in the structural accounting. For evaluation practice, a correspondence rate should
therefore be accompanied by the number of distinct rows, the actual choices
under alternative supported rules, and the action-copy control. These checks
separate an observed match from an opportunity to discriminate between rules;
they do not replace a test of prediction on unseen decisions.

\paragraph{Conditional outcomes and observed opportunities.}
Table~\ref{tab:followup-pairs} reports the two predeclared contrasts using only
observed nontechnical pairs. Astra has seven fully completed triplets and Sol
eight. L outcomes remain in the comparison; technical failures are reported
separately and are not converted into L. The directions do not support a
common improvement claim across these profiles, nor an equivalence or
no-effect conclusion. In Sol B, eight matching records coexist with two L
labels, illustrating why correspondence and continuation feasibility require
separate measurements. P certifies a conditional feasible continuation;
the experiments do not measure completed-task success.

% Generated from identity-bound existing follow-up sources.
\begin{table}[htbp]
\centering\small
\begin{tabular}{llrrrrrr}\toprule
Profile & Contrast & Pairs & Missing & L$\to$P & P$\to$L & Same & Net P\\\midrule
Astra & B--A & 7/8 & 1 & 1 & 0 & 6 & +1\\
Astra & C--B & 7/8 & 1 & 0 & 2 & 5 & -2\\
Sol & B--A & 8/8 & 0 & 0 & 1 & 7 & -1\\
Sol & C--B & 8/8 & 0 & 2 & 1 & 5 & +1\\
\bottomrule
\end{tabular}
\caption{Descriptive paired transitions run from A to B for B--A, and from B to C for C--B. Pairs are observed/planned; Missing counts technical missing pairs; Same includes P-to-P and L-to-L. Dividing Net P by observed pairs recovers the reported paired difference. No pooled estimate or population inference. C--B is the total condition contrast, not an isolated executor effect on the same fixed record.}
\label{tab:followup-pairs}
\end{table}


\paragraph{Where the outcome contrasts occur.}
The family breakdown in Table~\ref{tab:followup-families} locates every
observed P/L movement in the two continuous-control histories. All completed
route decisions were P, although each route menu contained an L option;
this pattern was not forced by having only P choices. The continuous menus
have 13 actions and the route menus three, so these purposeful contexts are
not a calibrated difficulty ladder. Appendix~\ref{app:followup-history}
reports every profile/history/arm cell. In particular, Sol's C--B net gain
of one P comprises two L-to-P and one P-to-L transitions. Net change alone
does not describe those opposing movements, and the underlying continuous
labels retain the numerical/runtime premises stated above.

% Generated from identity-bound existing follow-up sources.
\begin{table}[htbp]
\centering\small
\begin{tabular}{llrrrr}\toprule
Profile & Family & Assigned & Done & Tech. & P/L/U\\\midrule
Astra & Continuous & 12 & 12 & 0 & 9/3/0\\
Astra & Route & 12 & 10 & 2 & 10/0/0\\
Sol & Continuous & 12 & 12 & 0 & 8/4/0\\
Sol & Route & 12 & 12 & 0 & 12/0/0\\
\bottomrule
\end{tabular}
\caption{Follow-up outcomes by family, separately within each profile. Each family has two purposeful histories and two repeats per arm. Done denotes a completed decision, not full task success. P/L labels in the continuous family retain the numerical/runtime premises.}
\label{tab:followup-families}
\end{table}


\paragraph{A changed action need not change the task label.}
The follow-up contains two rule-sensitive records, both in C
(Table~\ref{tab:followup-rule-examples}). Mean and worst case select different
menu indices, while the bound label remains P in one and L in the other.
These are existing-record calculations, not new live counterfactual trials.
Together with the older panel's single label-changing example, they make
action sensitivity and task-label sensitivity separately inspectable.

% Generated from identity-bound existing follow-up sources.
\begin{table}[htbp]
\centering\small
\begin{tabular}{llrlll}\toprule
Profile & History & Repeat & Arm & Mean & Worst case\\\midrule
Astra & probe-01 & 2 & C & 1 (P) & 0 (P)\\
Sol & probe-02 & 1 & C & 11 (L) & 0 (L)\\
\bottomrule
\end{tabular}
\caption{The two rule-sensitive follow-up records. Entries give zero-based menu index and bound label. The action changes but the label does not; these are per-record calculations on existing menus, not live crossed decisions. Continuous labels remain conditional.}
\label{tab:followup-rule-examples}
\end{table}


Every completed C proposal already matched the record-selected action
(Astra 7/7, Sol 8/8). No observed C trial therefore tests the benefit of
overriding a conflicting proposal. The C--B contrast is not an isolated
mechanical correction effect on the same fixed record. Labels and
within-record rule diagnostics are tied to the existing menu contract;
they are not additional live crossing observations or independent evidence
about new task physics.

\section{Prospective pilot: does a frozen record predict a new choice?}
\label{sec:prospective}
We test whether a source decision record improves prediction of later model
choices beyond the public history and source action. Four new purposefully
constructed histories in the two disclosed task families are used separately
for Astra/max/Fast and Sol/max/Fast. Each source has two precommitted target
branches: eight target choices per profile, grouped in four histories.
The histories, records and target choices form a new panel. Each target
extends its source history by one observed transition under a fixed exogenous
action: continuous sources at step 6 branch to step 7 under $(0.6,0)$ or
$(-0.6,0)$; route sources at step 4 branch to step 5 by advancing one or three
nodes. These alternatives were set before source-model outputs and need not
follow the source proposal. Table~\ref{tab:prospective-branches} and
Appendix~\ref{app:prospective} give the construction and selection basis.

\paragraph{Information and output protocols.}
The source actor's final response contains an action, a numeric table with a
declared rule, and four brief fields for assumptions, uncertainty, criterion
and expected observation. The record is frozen before the source actor sees
either target. Two fresh-context observers of the same requested profile
predict the target actor's top-1 action with or without that record. Both
receive the source action, source and target histories, menus, and the
complete source and target instructions and output schemas. Only the record
differs between observer inputs. Their order is preallocated, allowances are
equal, and both predictions are locked before the fresh target actor runs.

The target's final response contains an action alone. Source and target may
both request up to three local analyses before their final response, and
their proposals are unforced. The target receives the public evidence and
source action, without the record or observer predictions. The comparison
therefore measures the record's added value for predicting an action under
this target protocol. Both observers predict the same target; the effect of
changing the source--target output protocol is not separately identified.

\begin{table}[htbp]
\centering\small
\begin{tabular}{llrrrrrl}\toprule
Panel & Profile & Pairs & Record & No record & Copy & Changed & Difference\\\midrule
Initial & Astra & 8/8 & 6/8 & 7/8 & 5/8 & 3/8 & -0.125\\
Initial & Sol & 7/8 & 3/7 & 3/7 & 0/7 & 7/7 & $[-0.125,+0.125]$\\
Combined & Astra & 8/8 & 6/8 & 7/8 & 5/8 & 3/8 & -0.125\\
Combined & Sol & 8/8 & 3/8 & 3/8 & 0/8 & 8/8 & 0.000\\
\bottomrule\end{tabular}
\caption{Prospective top-1 action prediction. Pairs retains the eight assigned
targets/profile; other fractions use available valid targets. Difference is
the equally weighted mean over four history clusters, with two target slots
each. The initial Sol interval is a missing-data identification bound over all
eight assigned targets, not a confidence interval; its observed 3/7 versus
3/7 comparison has difference zero. Copy uses the old source action; Changed
compares the target action with that action. Profiles are not ranked.}
\label{tab:prospective}
\end{table}

\begin{table}[htbp]
\centering\small
\begin{tabular}{lrrrr}\toprule
Profile & Both correct & Record only & No-record only & Both wrong\\\midrule
Astra & 5 & 1 & 2 & 0\\
Sol & 2 & 1 & 1 & 4\\
\bottomrule\end{tabular}
\caption{Joint top-1 prediction correctness for the combined panel, eight
targets per profile. Record only and No-record only identify which observer
is correct on the same target. Each cell uses one generation per role;
the counts describe paired outcomes, not case-level causal effects.}
\label{tab:prospective-paired}
\end{table}


\paragraph{Prediction results.}
The combined panel yields 6/8 versus 7/8 correct predictions for Astra and
3/8 versus 3/8 for Sol, with versus without the record. The mean of the four
within-history paired differences is $-0.125$ for Astra and $0$ for Sol.
Table~\ref{tab:prospective-paired} shows the joint outcomes: Sol's equal
totals include one target correctly predicted only with the record and
another correctly predicted only without it. Across these assignments,
no positive net accuracy gain is observed. Each cell uses one generation
per observer and target role; a discordant pair does not by itself identify
a case-level causal benefit or harm.

\paragraph{Where the prediction contrasts occur.}
Table~\ref{tab:prospective-families} adds a post hoc descriptive breakdown.
Sol's aggregate zero combines one additional correct prediction with the
record on continuous tasks and one fewer on route tasks. Astra's one fewer
correct prediction occurs in continuous control; both observers predict all
four route targets correctly. Each family contains two purposeful histories
with two linked targets each. The prespecified per-profile aggregate remains
primary; these counts locate the observed contrasts without estimating a
family-level causal effect. Table~\ref{tab:prospective-histories} gives all
eight history/profile units.

\begin{table}[htbp]
\centering\small
\begin{tabular}{llrrr}\toprule
Profile & Family & With record & Without record & Net correct\\\midrule
Astra & Continuous & 2/4 & 3/4 & $-1$\\
Astra & Route & 4/4 & 4/4 & $0$\\
Sol & Continuous & 3/4 & 2/4 & $+1$\\
Sol & Route & 0/4 & 1/4 & $-1$\\
\bottomrule\end{tabular}
\caption{Post hoc descriptive family breakdown of the combined prospective
panel. Each family has two purposeful histories with two linked targets each.
Net correct is the count with minus without the record. The prespecified
per-profile aggregate remains primary; these four-target contrasts do not
estimate family-level causal effects.}
\label{tab:prospective-families}
\end{table}


Copying the source action predicts 5/8 Astra and 0/8 Sol targets. The target
changes action relative to the source in 3/8 and 8/8 cases, respectively.
This baseline describes source--target action change; variability across
repeated runs of an identical target input was not measured.

\paragraph{Record structure and interpretation.}
All four source records per profile match the unforced source proposal.
Seven of the eight numeric tables contain a single row: four for Astra and
three for Sol. Agreement of the two supported rules is forced in these seven
tables and observed in the remaining two-distinct-row Sol table. No duplicate
row occurs. The semantic informativeness of the accompanying prose was not
independently evaluated. The treatment is the complete record, including its
prose; table row counts alone cannot explain the prediction result.

The reference controls verify protocol handling and prediction of supplied
hand-authored policies (Appendix~\ref{app:prospective}). Sensitivity to an
informative record in the scientific tasks remains a separate control target.
The observed contrast combines record information, observer use, the
source--target protocol change and one draw of each observer and target.
It does not separate those explanations for the net result. Task-matched
informative-record controls, a same-output-protocol comparison and repeated
identical-input decisions would distinguish them in a further study.
With four purposeful histories per profile and two linked targets per
history, the present pilot characterizes these assignments.

\paragraph{Execution segments.}
The original run produced fifteen targets. After inspecting its partial
results, we completed the remaining assigned Sol block in a separate segment
with the same source, inputs, profile and allowances and four fresh reference
checks. Table~\ref{tab:prospective} retains both panels. The completion adds
one target to an existing history; no completed scientific decision was
rerun. Appendix~\ref{app:prospective} retains the interrupted attempt, all
assigned roles and usage. The two segments may differ in time or provider
state; they form one descriptive pilot, not a confirmatory replication.

\section{Position relative to existing work}
\label{sec:related}
\paragraph{Execution and component diagnosis.}
\citet{lyu2023faithful} translate a problem into an intermediate representation
and use a deterministic solver to obtain an answer, locating their faithfulness
guarantee in the solution stage. Our C condition similarly binds execution
to an external record. We examine the additional evaluation questions raised
by constructing and comparing such records: selector-completion sensitivity,
correspondence without enforced execution, and prediction of later choices.
CAR estimates effects of interventions on agent steps through same-policy
downstream replay \citep{shah2026car}; AgentRx uses constraint evidence and a
judge for execution diagnosis \citep{barke2026agentrx}; and BLADE aligns evidence
from action stages for inspection \citep{guo2026blade}. RUU's component
diagnoses instead attach feasibility labels to actions selected from
external records. The diagnosed object and the outcome measure therefore
differ from stochastic-policy replay or failure attribution.

\paragraph{Alternative explanations and discriminating tests.}
\citet{jain2019attention} show that substantially different attention
distributions can preserve model predictions. \citet{wiegreffe2019attention}
argue that explanation claims depend on the model relationship being tested,
and introduce baselines, variance calibration and diagnostic tests to assess
attention's usefulness. Their debate distinguishes finding an alternative
from evaluating the explanatory value of the original. Our counterexample
concerns an analyst's external selector completion: the common best plan
on A2 makes the original-action match nondiscriminating, while A1 separates
the completed selectors. It identifies an ambiguity in that validation
criterion, rather than testing attention or establishing that the original
record is uninformative. This scope follows the broader requirement to
evaluate explanations against the system they describe
\citep{jacovi2020faithful}. Interventions on neural representations address
internal causal relations directly \citep{geiger2021causal}; our intervention
is on the external record and selection rule.

\paragraph{Prediction-based explanation evaluation.}
Leakage-Adjusted Simulatability (LAS) measures whether an explanation helps an
observer predict a model output while controlling for direct output leakage
\citep{hase2020las}. REV measures new label-relevant information beyond an
input-and-label baseline using conditional V-information \citep{chen2023rev}.
These approaches motivate assessing what an explanation adds to information
already available. \citet{madsen2024self} evaluate self-explanation consistency
through counterfactual, attribution and redaction tasks, while
\citet{mayne2026positive} test added explanation value for predicting a
reference model's answers on related inputs. They also use repeated
reference-model outputs to assess the prediction ceiling imposed by target
variability (their Section 4.3 and Appendix A.7); our single-generation pilot
does not estimate that ceiling.

Our prospective pilot applies the added-prediction question to records frozen
before precommitted target branches. Both observers already receive the source
action; neither sees the target choice before predicting it. The measured
quantity is the paired top-1 accuracy difference on later choices across
different source and target output protocols. LAS's leakage adjustment and
REV's information metric are distinct estimands. The pilot supplies a
source-auditable RUU case for this established evaluation question, alongside
record-structure and source-action-copy controls.

\paragraph{Connection to scientific discovery.}
BoxingGym evaluates experiment selection through expected information gain
and model discovery through predictions, including whether a scientist's
explanation helps another agent predict without further experiments
\citep{gandhi2025boxinggym}. RUU's present study concerns interpretation of
decision records and prediction of model choices. Prediction of withheld
environment observations and the utility of selected experiments are the
next behavioral targets in our broader research program. Keeping these
outcomes distinct connects the measurement lesson here to discovery
evaluation without treating action correspondence as discovery success.

\section{Interpretation and next research questions}
The case study yields a practical sequence for evaluating decision records.
First, establish which records support the proposed operation. The transfer
panel's zero eligible pairs locate an applicability bottleneck. Second,
specify the analyst's completion and test whether the validation cases
separate plausible alternatives. The pilot's common best plan shows why a
successful original-action match can leave this choice open. Third, report
correspondence together with structural opportunities for rule discrimination.
The unforced-proposal follow-up demonstrates their coexistence: all sixteen
required records match, while fourteen have mathematically forced rule
agreement. Finally, evaluate added prediction on new decisions. The prospective
pilot supplies that separate measurement and observes no positive net accuracy
gain in its assigned cases.

These checks also determine the appropriate unit of evidence. Under the fixed
numeric contract, a per-record rule scan exposes the sensitivity in the paired
display (Appendix~\ref{app:structural-control}). Action changes and P/L/U label
changes answer different questions, as the two rule-sensitive but
label-preserving records illustrate. The proposed audit makes these
distinctions inspectable; its effect on downstream diagnostic accuracy remains
to be evaluated against independently established diagnoses.

\paragraph{Scope of the evidence.}
The histories are purposefully selected within two task families. Continuous
P/L labels retain the numerical and historical-runtime premises in
Appendix~\ref{app:runtime-proof}; P certifies a feasible continuation, not a
completed task. All completed C proposals already agree with their selectors,
so correcting a conflicting proposal has no observed denominator. The older
panel's seven analysis errors, including its sole label-changing
rule-sensitive record, have no isolated behavioral-effect estimate. Record
correspondence and external substitutions address observed outputs; internal
causal faithfulness requires a different intervention. In the prospective
pilot, same-protocol targets, repeated identical-input targets, and a control
that demonstrates use of informative records in the scientific tasks would
help distinguish protocol transfer, target variability and observer sensitivity.

\paragraph{Research under uncertainty.}
Three questions guide the next stage: how a stated assumption relates to the
next experiment choice; whether that experiment resolves more uncertainty
than an equally resourced valid fixed plan; and whether a revised explanation
improves predictions for cases withheld during its construction. Each question
calls for its own scored behavioral outcome. The present prospective pilot
measures model-choice prediction. Prediction of new world observations and
the utility of adaptive experimentation will extend this program from
interpreting decision records to evaluating scientific behavior.

\section{Reader access and reproducibility}
\label{sec:reader-access}
The anonymous supplement supplies the manuscript sources, current-study
mechanisms, visible interaction evidence and source-bound readers.
Appendix~\ref{app:reader-access} lists the material and check boundaries for
each evidence set. One archive contains both the main and prospective
components in a lossless text representation; its reader restores their
exact bytes, verifies their manifests and runs the evidence checks.
From the extracted supplementary archive, run:
\begin{verbatim}
python3.12 -I -B -S verify.py
\end{verbatim}
The command recomputes the reported source-level results, checks the paired
prediction outcomes and selector-discrimination example, and runs the
existing numerical replays. No model account is needed. The README also
provides permanent extraction and corruption-test commands. Source
recalculation and present-host numerical replay establish reproducibility
of the supplied evidence; historical/provider attestation and independent
laboratory replication remain separate verification levels.

\paragraph{Responsible interpretation.}
The experiments use simulated environments and model outputs. The checks
evaluate external decision records; assessing internal reasoning, intentions
or deployment safety requires evidence for those distinct claims. The
disclosed mechanisms support inspection of this fixed study. Future
benchmark mechanisms have a separate disclosure policy.
